\documentclass[10pt,a4paper]{article}

\usepackage[margin=2cm]{geometry}
\usepackage[T1]{fontenc}
\usepackage{lmodern}
\usepackage{amsmath}
\usepackage{booktabs}
\usepackage{tabularx}
\usepackage{enumitem}
\usepackage[hidelinks]{hyperref}

\setlist{nosep, leftmargin=*}
\setlength{\parskip}{3pt}
\setlength{\parindent}{0pt}

\title{\vspace{-1.6cm}\textbf{Project Proposal}\\[2pt]
\large CNN vs Vision Transformer for Image Classification:\\
Inductive Bias, Dataset Size, Augmentation and Compute Cost}
\author{Machine Learning and Applications --- Topic 7\\
[Member 1 name, student ID] (team leader)\quad [Member 2 name, student ID]}
\date{}

\begin{document}
\maketitle
\vspace{-1.2cm}

\section{Problem and Research Questions}
Convolutional neural networks (CNNs) build in strong assumptions about images: locality and
weight sharing, which give translation equivariance~\cite{he2016resnet}. Vision Transformers
(ViTs) treat an image as a sequence of patches and rely on global self-attention, with far weaker
built-in assumptions~\cite{dosovitskiy2021vit}. Dosovitskiy et al.\ report that ViTs trail
comparable ResNets when pretrained on ImageNet-1k alone but match or exceed them with larger
pretraining datasets, suggesting that weak inductive bias must be compensated by data.
Later work shows that the gap also depends on training recipe and augmentation~\cite{touvron2021deit,
steiner2022howtotrain}, that modernised CNNs can match ViTs~\cite{liu2022convnext}, and that the
two families learn different internal representations~\cite{raghu2021vitsee}.
We study these claims in a controlled, small-scale setting on CIFAR-10~\cite{krizhevsky2009cifar}:
\begin{enumerate}
  \item[\textbf{RQ1}] \textbf{Dataset size:} how does accuracy change as the training set shrinks
  from 100\% to 1\%, for pretrained and from-scratch CNNs and ViTs?
  \item[\textbf{RQ2}] \textbf{Augmentation:} which family benefits more from stronger augmentation?
  \item[\textbf{RQ3}] \textbf{Compute cost:} how do the models compare in accuracy per training
  time, GPU memory and number of parameters?
  \item[\textbf{RQ4}] \textbf{Interpretation:} do Grad-CAM~\cite{selvaraju2017gradcam} and attention
  rollout~\cite{abnar2020rollout} maps show different evidence used by each family?
\end{enumerate}
\textbf{Hypothesis.} Trained from scratch, CNNs should outperform a ViT of similar size on small
training fractions, and the gap should narrow as data grows. Pretraining should reduce this gap,
and ViTs should gain more from strong augmentation.

\section{Main Papers and Methods}
We reproduce the core comparison of ViT~\cite{dosovitskiy2021vit} against ResNet~\cite{he2016resnet}
and ConvNeXt~\cite{liu2022convnext}. ResNet learns residual mappings
$y = \mathcal{F}(x,\{W_i\}) + x$; ViT applies self-attention
$\mathrm{softmax}(QK^\top/\sqrt{d_k})V$ over patch embeddings plus learned position embeddings.
Pretrained models come from timm~\cite{wightman2019timm}, using ImageNet-1k weights for all models
(ResNet-50 trained with the recipe of~\cite{wightman2021resnetstrikes}, ViT-S/16
from~\cite{steiner2022howtotrain}) so that pretraining data is the same across families.

\begin{tabularx}{\linewidth}{@{}lllXr@{}}
\toprule
\textbf{Track} & \textbf{Model} & \textbf{Weights} & \textbf{Input} & \textbf{Params} \\
\midrule
Fine-tune & ResNet-50 & ImageNet-1k & $224\times224$ (upscaled) & 23.5M \\
Fine-tune & ConvNeXt-Tiny & ImageNet-1k & $224\times224$ (upscaled) & 27.8M \\
Fine-tune & ViT-S/16 & ImageNet-1k & $224\times224$ (upscaled) & 21.7M \\
Scratch & SmallCNN (own) & random & $32\times32$ (native) & TBD \\
Scratch & SmallViT (own, $4\times4$ patches) & random & $32\times32$ (native) & matched to SmallCNN \\
\bottomrule
\end{tabularx}
\vspace{2pt}

\textbf{Baseline.} Logistic regression on raw pixels checks the data pipeline and gives a reference.

\textbf{Training.} AdamW~\cite{loshchilov2019adamw} (weight decay 0.05), one warmup epoch then cosine
decay~\cite{loshchilov2017sgdr}, batch size 32 and mixed precision. Fine-tuning starts from learning
rates $10^{-3}$ (new head) and $10^{-4}$ (backbone) for 8 epochs; scratch models use $10^{-3}$ for
60 epochs. Every model receives the same tuning budget on the validation set, and the search space
will be reported. Augmentation levels: none; basic (random crop, horizontal flip); strong
(adding RandAugment~\cite{cubuk2020randaugment}).

\section{Data, Experiments and Metrics}
\textbf{Data.} CIFAR-10~\cite{krizhevsky2009cifar}: 50,000 training and 10,000 test images in 10
classes. A stratified 10\% of the training set is held out for validation; the test set is used
once per model, after model selection on validation.

\textbf{Experiments.}
\begin{itemize}
  \item \textbf{Learning curves (RQ1):} training fractions 1, 5, 10, 25, 50, 100\%, using
  stratified, nested subsets; 5 models $\times$ 6 fractions $\times$ 3 seeds.
  \item \textbf{Augmentation (RQ2):} none / basic / strong at fixed training fractions.
  \item \textbf{Ablations:} pretrained vs random initialisation (ResNet-50, ViT-S/16); frozen
  backbone (linear probe) vs full fine-tuning.
  \item \textbf{Interpretation (RQ4):} Grad-CAM for CNNs and attention rollout for ViTs, on one
  shared set of correctly and incorrectly classified test images.
  \item \textbf{Extension (if time allows):} robustness on CIFAR-10-C~\cite{hendrycks2019robustness}
  and shape/texture bias~\cite{geirhos2019texture}.
\end{itemize}

\textbf{Metrics.} Test accuracy and macro-F1 as mean $\pm$ std over seeds 0, 1, 2; training time;
peak GPU memory; parameter count.

\section{Task Allocation}
\begin{tabularx}{\linewidth}{@{}lX@{}}
\toprule
\textbf{Member} & \textbf{Responsibilities} \\
\midrule
[Member 1] & CNN track: ResNet-50 and ConvNeXt-Tiny fine-tuning, SmallCNN, Grad-CAM; shared data
pipeline and training loop. \\
[Member 2] & ViT track: ViT-S/16 fine-tuning, SmallViT, attention rollout. \\
Both & Logistic regression baseline, literature review, result analysis, report and presentation. \\
\bottomrule
\end{tabularx}

\section{Compute Resources and Timeline}
Development uses a laptop with an NVIDIA RTX 3050 Ti (4 GB), PyTorch 2.11 and timm 1.0.31. A
preliminary benchmark at $224\times224$ with batch size 32 measured 174 (ResNet-50), 134
(ConvNeXt-Tiny) and 247 (ViT-S/16) training images per second, about 3--6 minutes per epoch on
the full training split. The full learning-curve grid is estimated at roughly 3--4 GPU hours per
pretrained model, before tuning; Kaggle or Colab GPUs serve as a fallback, and the hardware will be
recorded for every run. Code is versioned on GitHub.

\textbf{Timeline.} Week 1: proposal. Weeks 2--3: literature review, data pipeline, baseline
(progress report). Weeks 4--6: fine-tuning, scratch models, learning curves, ablations.
Weeks 7--8: interpretation, report and presentation.

\textbf{Expected limitations.} Upscaled $32\times32$ images carry little fine detail; pretrained
and scratch models use different resolutions; and our small scale cannot reproduce the large
pretraining regime of~\cite{dosovitskiy2021vit}.

\small
\begin{thebibliography}{99}
\setlength{\itemsep}{0pt}
\bibitem{he2016resnet} K.~He, X.~Zhang, S.~Ren, J.~Sun. Deep Residual Learning for Image
Recognition. \emph{CVPR}, 2016. arXiv:1512.03385.
\bibitem{dosovitskiy2021vit} A.~Dosovitskiy et al. An Image is Worth 16x16 Words: Transformers for
Image Recognition at Scale. \emph{ICLR}, 2021. arXiv:2010.11929.
\bibitem{liu2022convnext} Z.~Liu, H.~Mao, C.-Y.~Wu, C.~Feichtenhofer, T.~Darrell, S.~Xie. A ConvNet
for the 2020s. \emph{CVPR}, 2022. arXiv:2201.03545.
\bibitem{touvron2021deit} H.~Touvron et al. Training Data-Efficient Image Transformers \&
Distillation through Attention. \emph{ICML}, 2021. arXiv:2012.12877.
\bibitem{steiner2022howtotrain} A.~Steiner et al. How to Train Your ViT? Data, Augmentation, and
Regularization in Vision Transformers. \emph{TMLR}, 2022. arXiv:2106.10270.
\bibitem{raghu2021vitsee} M.~Raghu, T.~Unterthiner, S.~Kornblith, C.~Zhang, A.~Dosovitskiy. Do Vision
Transformers See Like Convolutional Neural Networks? \emph{NeurIPS}, 2021. arXiv:2108.08810.
\bibitem{krizhevsky2009cifar} A.~Krizhevsky. Learning Multiple Layers of Features from Tiny Images.
Technical report, University of Toronto, 2009.
\bibitem{selvaraju2017gradcam} R.~R.~Selvaraju et al. Grad-CAM: Visual Explanations from Deep
Networks via Gradient-Based Localization. \emph{ICCV}, 2017. arXiv:1610.02391.
\bibitem{abnar2020rollout} S.~Abnar, W.~Zuidema. Quantifying Attention Flow in Transformers.
\emph{ACL}, 2020. arXiv:2005.00928.
\bibitem{wightman2019timm} R.~Wightman. PyTorch Image Models.
\url{https://github.com/huggingface/pytorch-image-models}, 2019.
\bibitem{wightman2021resnetstrikes} R.~Wightman, H.~Touvron, H.~J\'egou. ResNet Strikes Back: An
Improved Training Procedure in timm. arXiv:2110.00476, 2021.
\bibitem{loshchilov2019adamw} I.~Loshchilov, F.~Hutter. Decoupled Weight Decay Regularization.
\emph{ICLR}, 2019. arXiv:1711.05101.
\bibitem{loshchilov2017sgdr} I.~Loshchilov, F.~Hutter. SGDR: Stochastic Gradient Descent with Warm
Restarts. \emph{ICLR}, 2017. arXiv:1608.03983.
\bibitem{cubuk2020randaugment} E.~D.~Cubuk, B.~Zoph, J.~Shlens, Q.~V.~Le. RandAugment: Practical
Automated Data Augmentation with a Reduced Search Space. \emph{CVPR Workshops}, 2020.
arXiv:1909.13719.
\bibitem{hendrycks2019robustness} D.~Hendrycks, T.~Dietterich. Benchmarking Neural Network Robustness
to Common Corruptions and Perturbations. \emph{ICLR}, 2019. arXiv:1903.12261.
\bibitem{geirhos2019texture} R.~Geirhos et al. ImageNet-Trained CNNs Are Biased towards Texture;
Increasing Shape Bias Improves Accuracy and Robustness. \emph{ICLR}, 2019. arXiv:1811.12231.
\end{thebibliography}

\end{document}
